This appendix collects a few explicit computations that support Sections~\ref{subsec:geompropagation} and~\ref{sec:global_torsion_acs_narrative_enhanced_lc}. The guiding principle is the same as in the main text: only calculations that clarify already-established structures are recorded here. No new program is introduced.

\subsection{Horocyclic coordinates and the arithmetic grid in \texorpdfstring{$\mathfrak{E}_0$}{E0}}
\label{appC:horocyclic-grid}

Recall that on the upper half-plane
\[
\mathcal H=\{(x,y)\in\mathbb R^2\mid y>0\}
\]
we consider the metric and assignment field
\[
 g_{\mu,\lambda}=\frac1{y^2}\left(\frac{dx^2}{\mu^2}+\frac{dy^2}{\lambda^2}\right),
 \qquad
 a(x,y)=-\frac{x}{y}.
\]
Set
\[
 u=x,
 \qquad
 v=\lambda^{-1}\log y.
\]
Then $y=e^{\lambda v}$ and therefore
\[
 dx=du,
 \qquad
 dy=\lambda e^{\lambda v}\,dv.
\]
Substituting into the metric gives
\[
 g_{\mu,\lambda}
 =\frac{1}{e^{2\lambda v}}\left(\frac{du^2}{\mu^2}+\frac{\lambda^2e^{2\lambda v}dv^2}{\lambda^2}\right)
 =e^{-2\lambda v}\frac{du^2}{\mu^2}+dv^2,
\]
while the assignment becomes
\[
 a(u,v)=-\frac{u}{e^{\lambda v}}=-u e^{-\lambda v}.
\]
This is the horocyclic form used in Section~\ref{subsec:geompropagation}.

The elementary arithmetic moves can now be read off directly.
For an additive step $s\in\mathbb R$ and a multiplicative factor $k>0$, define
\[
 X_s(u,v)=(u-se^{\lambda v},v),
 \qquad
 Y_k(u,v)=\bigl(u,v-\lambda^{-1}\log k\bigr).
\]
Then
\[
 a\bigl(X_s(u,v)\bigr)
 =-\bigl(u-se^{\lambda v}\bigr)e^{-\lambda v}
 =-ue^{-\lambda v}+s
 =a(u,v)+s,
\]
and
\[
 a\bigl(Y_k(u,v)\bigr)
 =-u e^{-\lambda(v-\lambda^{-1}\log k)}
 =-u e^{-\lambda v}e^{\log k}
 =k a(u,v).
\]
So $X_s$ realizes addition and $Y_k$ realizes multiplication.

\begin{proposition}
\label{prop:appC-one-cell-torsion}
For one additive step $s$ and one multiplicative factor $k=e^{\Lambda}$, the two elementary histories
\[
\gamma_{\square}=(\oplus_s,\otimes_{\Lambda}),
\qquad
\bar\gamma_{\square}=(\otimes_{\Lambda},\oplus_s)
\]
have torsion
\[
\tau(\gamma_{\square})=s(e^{\Lambda}-1)=s(k-1).
\]
Moreover, this torsion is exactly the weighted area of the rectangle
$[0,s]\times[0,\Lambda]$ in the ACS:
\[
\tau(\gamma_{\square})
 =\int_0^s\int_0^{\Lambda} e^M\,dM\,dA.
\]
\end{proposition}

\begin{proof}
Direct evaluation gives
\[
\nu_x(\gamma_{\square})=e^{\Lambda}(x+s),
\qquad
\nu_x(\bar\gamma_{\square})=e^{\Lambda}x+s,
\]
so
\[
\tau(\gamma_{\square})
 =e^{\Lambda}(x+s)-(e^{\Lambda}x+s)
 =s(e^{\Lambda}-1).
\]
On the other hand,
\[
\int_0^s\int_0^{\Lambda} e^M\,dM\,dA
 =\int_0^s (e^{\Lambda}-1)\,dA
 =s(e^{\Lambda}-1).
\]
The two expressions coincide.
\end{proof}

Thus the finite torsion carried by a single arithmetic cell is already the exact weighted-area quantity that later appears globally in the ACS. The infinitesimal law
\[
 d\tau = \mu\lambda\,du\,dv
\]
from Section~\ref{sec:flowequation} is simply the first-order expansion of this finite formula in coordinates adapted to the grid.

\begin{corollary}
\label{cor:appC-rectangular-family}
Let
\[
\gamma_{m,n}=\bigl(\oplus_1\bigr)^m\bigl(\otimes_{\log 2}\bigr)^n,
\qquad
\bar\gamma_{m,n}=\bigl(\otimes_{\log 2}\bigr)^n\bigl(\oplus_1\bigr)^m.
\]
Then
\[
\tau(\gamma_{m,n})=m(2^n-1).
\]
In particular, taking $m=n$ recovers the sequence $1,6,21,\dots$ from Section~\ref{subsec:gridsandtorsion}.
\end{corollary}

\begin{proof}
We have
\[
\nu_x(\gamma_{m,n})=2^n(x+m),
\qquad
\nu_x(\bar\gamma_{m,n})=2^n x+m.
\]
Hence
\[
\tau(\gamma_{m,n})=2^n(x+m)-(2^n x+m)=m(2^n-1).
\]
\end{proof}

\subsection{The radial formula in \texorpdfstring{$\mathfrak{E}_1$}{E1}}
\label{appC:first-kind-radial}

Section~\ref{subsec:firstkindspace} records the first kind space $\mathfrak{E}_1$ in the Poincar\'e disc with a single isolated zero. The assignment formula there is a direct consequence of the hyperbolic distance function.

\begin{lemma}
\label{lem:appC-radial-formula}
Let $r\in(0,1)$ be the Euclidean radial coordinate on the Poincar\'e disc with curvature $-\lambda^2$, and let
\[
 s(r)=\frac{2}{\lambda}\operatorname{arctanh}(r)
\]
be the hyperbolic distance from the center. If the radial arithmetic flow starts from the singular zero and satisfies
\[
 a(s)=\frac{\mu}{\lambda}\sinh(\lambda s),
\]
then
\[
 a(r)=\frac{\mu}{\lambda}\frac{2r}{1-r^2}.
\]
\end{lemma}

\begin{proof}
From
\[
 s(r)=\frac{2}{\lambda}\operatorname{arctanh}(r)
\]
we obtain
\[
 \lambda s = 2\operatorname{arctanh}(r).
\]
Hence
\[
 e^{\lambda s}=\frac{1+r}{1-r},
 \qquad
 e^{-\lambda s}=\frac{1-r}{1+r}.
\]
Therefore
\[
 \sinh(\lambda s)
 =\frac{e^{\lambda s}-e^{-\lambda s}}{2}
 =\frac12\left(\frac{1+r}{1-r}-\frac{1-r}{1+r}\right)
 =\frac12\cdot\frac{(1+r)^2-(1-r)^2}{1-r^2}
 =\frac{2r}{1-r^2}.
\]
Substituting this into $a(s)=\frac{\mu}{\lambda}\sinh(\lambda s)$ gives the claimed formula.
\end{proof}

Two asymptotic regimes are immediate:
\[
 a(r)\sim \frac{2\mu}{\lambda}r
 \quad (r\to 0),
 \qquad
 a(r)\sim \frac{\mu}{\lambda}\frac{1}{1-r}
 \quad (r\to 1^-).
\]
So the singular center behaves linearly to first order, while the assignment diverges at the ideal boundary.

\subsection{An inductive derivation of the ACS evaluation formula}
\label{appC:acs-induction}

The proof in Section~\ref{sec:global_torsion_acs_narrative_enhanced_lc} explains the evaluation formula by reading off the future multiplicative charge of each additive edge. For later reference it is convenient to record the same statement in an inductive form.

For a prefix
\[
\gamma^{(k)}=(g_1,\dots,g_k)
\]
of an arithmetic history, write
\[
 M_k:=\sum_{\substack{1\le j\le k\\ g_j=\otimes_{\lambda_j}}}\lambda_j
\]
for its total multiplicative charge. If $g_j=\oplus_{\mu_j}$ with $j\le k$, define the future multiplicative charge inside the prefix by
\[
 M^+_{j,k}:=\sum_{\substack{j<\ell\le k\\ g_\ell=\otimes_{\lambda_\ell}}}\lambda_\ell.
\]

\begin{proposition}
\label{prop:appC-prefix-evaluation}
For every prefix $\gamma^{(k)}$ and every initial value $x$,
\[
\nu_x\bigl(\gamma^{(k)}\bigr)
 = e^{M_k}x
   + \sum_{\substack{1\le j\le k\\ g_j=\oplus_{\mu_j}}}
     \mu_j e^{M^+_{j,k}}.
\]
\end{proposition}

\begin{proof}
We argue by induction on $k$.

For $k=0$ the history is empty, $M_0=0$, and the formula reads $\nu_x(\varnothing)=x$, which is true.

Assume the formula holds for $k$ and consider $\gamma^{(k+1)}$.
There are two cases.

\smallskip
\noindent
\emph{Case 1:} $g_{k+1}=\oplus_{\mu_{k+1}}$.
Then $M_{k+1}=M_k$, and
\[
\nu_x\bigl(\gamma^{(k+1)}\bigr)
 =\nu_x\bigl(\gamma^{(k)}\bigr)+\mu_{k+1}.
\]
By the induction hypothesis,
\[
\nu_x\bigl(\gamma^{(k+1)}\bigr)
 =e^{M_k}x
 +\sum_{\substack{1\le j\le k\\ g_j=\oplus_{\mu_j}}}\mu_j e^{M^+_{j,k}}
 +\mu_{k+1}.
\]
Since the new step is additive, the future multiplicative charges of all earlier additive steps remain unchanged, and the new additive term has future multiplicative charge $0$. Thus the required formula holds for $k+1$.

\smallskip
\noindent
\emph{Case 2:} $g_{k+1}=\otimes_{\lambda_{k+1}}$.
Then
\[
\nu_x\bigl(\gamma^{(k+1)}\bigr)
 =e^{\lambda_{k+1}}\nu_x\bigl(\gamma^{(k)}\bigr).
\]
Applying the induction hypothesis gives
\[
\nu_x\bigl(\gamma^{(k+1)}\bigr)
 =e^{\lambda_{k+1}}e^{M_k}x
 +\sum_{\substack{1\le j\le k\\ g_j=\oplus_{\mu_j}}}
   \mu_j e^{\lambda_{k+1}}e^{M^+_{j,k}}.
\]
Now $M_{k+1}=M_k+\lambda_{k+1}$, and every earlier additive step acquires one extra future multiplicative factor $e^{\lambda_{k+1}}$. Hence
\[
 e^{\lambda_{k+1}}e^{M^+_{j,k}} = e^{M^+_{j,k+1}},
\]
which yields the desired formula.
\end{proof}

Taking $k=n$ gives the evaluation formula from the main text. Reading the reversed history $\bar\gamma$ instead of $\gamma$ turns the same sum into a line integral along the reversed ACS path:
\[
\nu_x(\gamma)=e^{M_\gamma}x+\int_{C_{\bar\gamma}}e^M\,dA.
\]
The integral notation is simply a compact way to encode the same weighted sum.

\subsection{Explicit torsion formulas in the ACS}
\label{appC:explicit-torsion-formulas}

Let
\[
\gamma=(g_1,\dots,g_n)
\]
be an arithmetic history. For each additive step $g_j=\oplus_{\mu_j}$, define both the future and past multiplicative charges by
\[
 M_j^+ := \sum_{\substack{k>j\\ g_k=\otimes_{\lambda_k}}}\lambda_k,
 \qquad
 M_j^- := \sum_{\substack{k<j\\ g_k=\otimes_{\lambda_k}}}\lambda_k.
\]
Then the forward and reversed evaluations read
\[
\nu_x(\gamma)=e^{M_\gamma}x+\sum_{j\in I_A}\mu_j e^{M_j^+},
\qquad
\nu_x(\bar\gamma)=e^{M_\gamma}x+\sum_{j\in I_A}\mu_j e^{M_j^-},
\]
where $I_A$ denotes the set of additive indices.

\begin{corollary}
\label{cor:appC-explicit-torsion-formula}
For every arithmetic history $\gamma$,
\[
\tau(\gamma)
 =\sum_{j\in I_A}\mu_j\bigl(e^{M_j^+}-e^{M_j^-}\bigr).
\]
\end{corollary}

\begin{proof}
Subtract the two formulas above. The common term $e^{M_\gamma}x$ cancels.
\end{proof}

This expression is often the fastest way to compute torsion in concrete examples. It also makes two structural facts transparent.
First, only additive steps contribute directly to the sum; multiplicative steps act by reweighting those contributions. Second, a step contributes nothing precisely when its total multiplicative charge before and after are the same.

\begin{corollary}
\label{cor:appC-block-paths}
Suppose $\gamma$ is a block path of the form
\[
\gamma=(\oplus_{\mu_1},\dots,\oplus_{\mu_r},\otimes_{\lambda_1},\dots,\otimes_{\lambda_s}),
\]
and let
\[
\Lambda:=\lambda_1+\cdots+\lambda_s.
\]
Then
\[
\tau(\gamma)=\bigl(e^{\Lambda}-1\bigr)(\mu_1+\cdots+\mu_r).
\]
\end{corollary}

\begin{proof}
For every additive step in such a block path one has $M_j^+=\Lambda$ and $M_j^-=0$. The previous corollary therefore gives
\[
\tau(\gamma)
 =\sum_{j=1}^r \mu_j(e^{\Lambda}-1)
 =\bigl(e^{\Lambda}-1\bigr)(\mu_1+\cdots+\mu_r).
\]
\end{proof}

\paragraph{A signed example.}
To illustrate the oriented nature of the ACS filling, consider
\[
\gamma=(\oplus_1,\otimes_{\log 2},\oplus_{-1}).
\]
Then
\[
\nu_x(\gamma)=2(x+1)-1=2x+1,
\]
whereas
\[
\nu_x(\bar\gamma)=(x-1)\cdot 2+1=2x-1.
\]
So
\[
\tau(\gamma)=2.
\]
In ACS language,
\[
C_\gamma:(0,0)\to(1,0)\to(1,\log 2)\to(0,\log 2),
\]
while
\[
C_{\bar\gamma}:(0,0)\to(-1,0)\to(-1,\log 2)\to(0,\log 2).
\]
These two paths together bound the rectangle
\[
[-1,1]\times[0,\log 2],
\]
so the triple identity becomes
\[
\tau(\gamma)
 =\int_{-1}^1\int_0^{\log 2} e^M\,dM\,dA
 =2.
\]
This is a useful reminder that the ACS formalism is not restricted to monotone positive staircases: the same weighted-area language also handles signed additive data through oriented chains.
